% id: 199-33
% book: 베이직쎈 대수 · 12 수학적 귀납법
% section: 개념 70 수학적 귀납법
% passage: 다음 ㈎, ㈏에 알맞은 것을 구하시오.
% figure: none
% answer: 
% answer_source: 
% variant_level: 0 (원본 전사)
% 이 파일은 build.mjs 가 items.json 에서 만든다. 고칠 때는 items.json 을 고친다.
\smash{\raisebox{1.5mm}{\dmkichul{베이직쎈 대수 199쪽 33번}}}
자연수 $n$에 대한 명제 $p(n)$이 모든 자연수 $n$에 대하여 성립함을 수학적 귀납법을 이용하여 증명하려면 다음 두 가지를 보이면 된다.\exprbox{\begin{minipage}{0.96\linewidth}\raggedright\lineskip=4pt (i) \fbox{㈎}일 때, 명제 $p(n)$이 성립한다.\\ (ii) $n=k$일 때 명제 $p(n)$이 성립한다고 가정하면\\ \hspace*{1.5em}$n=$\fbox{㈏}일 때도 명제 $p(n)$이 성립한다.\end{minipage}}
